%% ============================================================
%%  AttnFuse — Speaker Notes for the Final Project Presentation
%%  One section per slide in AttnFuse_slides.pptx
%%  Written in plain, beginner-friendly language.
%% ============================================================
\documentclass[11pt]{article}

\usepackage[margin=1in]{geometry}
\usepackage[T1]{fontenc}
\usepackage{lmodern}
\usepackage{microtype}
\usepackage{amsmath,amssymb}
\usepackage{graphicx}
\usepackage{booktabs}
\usepackage{listings}
\usepackage{xcolor}
\usepackage[colorlinks=true,linkcolor=blue!60!black,
            citecolor=blue!60!black,urlcolor=blue!60!black]{hyperref}
\usepackage{parskip}
\usepackage{enumitem}
\usepackage{titlesec}

\titleformat{\section}{\Large\bfseries\color{blue!60!black}}{Slide \thesection.}{0.6em}{}

\definecolor{codebg}{HTML}{F4F4F4}
\definecolor{tipbg}{HTML}{FFF5E1}
\lstset{
  language=Python,
  backgroundcolor=\color{codebg},
  basicstyle=\ttfamily\small,
  frame=single,
  rulecolor=\color{black!20},
  breaklines=true,
  showstringspaces=false,
  tabsize=2,
}

\newcommand{\tip}[1]{%
  \par\smallskip
  \noindent\fcolorbox{orange!50}{tipbg}{%
    \parbox{0.96\linewidth}{\textbf{What to say out loud:} #1}}%
  \par\smallskip
}

\title{\textbf{AttnFuse — Speaker Notes}\\[4pt]
       \large Final Project Presentation, CS 790/657}
\author{Varun Kumar Dasoju}
\date{}

\begin{document}
\maketitle

\noindent\textit{This document is your script. Each section maps one-to-one to a slide
in \texttt{AttnFuse\_slides.pptx}. The orange boxes are exactly what to say out
loud; the text above them gives you background you can paraphrase if asked.
The whole talk runs comfortably in 10--12 minutes.}

%% =============================================================
\section{Title — ``AttnFuse''}
%% =============================================================

\tip{Hi everyone. My project is called AttnFuse. It is a small programming
language for attention, the most expensive part of modern AI models like ChatGPT.
You write the attention as a short Python recipe, and AttnFuse turns it into
a fast GPU program automatically. Over the next ten minutes I will tell you
why this matters, how it works, and how fast it ended up being.}

%% =============================================================
\section{What Is This Project About?}
%% =============================================================

The audience may have no idea what ``attention'' means yet. Don't define it
here. Just give them the one-sentence pitch and a reason to care.

\tip{Big AI models read text the way you read a sentence: they figure out
which words matter most. The piece of the model that does that is called
\textit{attention}, and it is the slowest part of the whole system. Today,
if a researcher wants to invent a new flavour of attention --- say, one that
only looks at the last 256 words --- they have to write hundreds of lines
of low-level GPU code, which takes weeks. AttnFuse is a small language that
lets you describe the new flavour in five lines of Python, and it spits out
the fast GPU code for you automatically.}

%% =============================================================
\section{What Is ``Attention'' in AI? — A Simple Picture}
%% =============================================================

Use the highlighter analogy. Don't get into math yet.

\tip{Imagine you are reading the sentence ``She bought a red apple.''
When you reach the word ``apple,'' your brain glances back at ``red'' to
understand which kind of apple. An AI model does exactly the same thing:
for every word it is currently reading, it glances back at all the other
words and gives each one a little score that says ``how relevant are
you to me right now?'' Those scores form a big grid --- one row per word
you are reading, one column per word you are looking at. The model then
blends all the words together, weighted by those scores. That blending
is what we call \textit{attention}.}

%% =============================================================
\section{The Math, Stripped Down}
%% =============================================================

This is the only math slide. Keep it brief; don't blackboard derivations.

\tip{Under the hood, every word becomes three short lists of numbers.
Think of them as a question, a label, and a piece of info that the word
carries. The first matrix multiply --- Q times K-transpose --- asks every
question against every label, and gives back a giant scoreboard. Softmax
turns the raw scores into percentages that add up to one. Then a second
matrix multiply uses those percentages to mix the pieces of info together.
That's the whole recipe. Every fancy variant of attention --- causal,
sliding-window, ALiBi, RoPE --- is just a small tweak on this scaffold.}

%% =============================================================
\section{Why Is Attention Slow?}
%% =============================================================

Drive home the quadratic cost. People who have heard the words ``context
length'' should now understand why scaling it up is so painful.

\tip{Here is the catch. If your text is N words long, the scoreboard has
N times N entries. At 4096 words --- a fairly short article --- that's
already sixteen million entries per attention head, per layer. A real
model has dozens of heads and dozens of layers, so you end up with tens
of billions of entries. The GPU is fast at math, but it is much slower at
moving data in and out of memory. Writing that giant scoreboard out and
reading it back in is what kills the speed. Standard PyTorch makes it
worse by launching five separate GPU programs in a row --- one for the
matmul, one for the scaling, one for the mask, one for softmax, one for
the final blend. AttnFuse fuses all of those into one.}

%% =============================================================
\section{FlashAttention: The Hand-Optimised Hero}
%% =============================================================

Acknowledge prior work clearly. This is what we are inspired by but not
just copying.

\tip{In 2022, Tri Dao published FlashAttention. The clever idea was: never
write the full scoreboard. Instead, process the input in small tiles, keep
a running max and a running sum on the GPU's tiny fast memory, and update
them as each tile arrives. The math comes out exactly the same, but the
big slow scoreboard never gets written down. FlashAttention is two to four
times faster than naive PyTorch, and it uses a tiny fraction of the memory.
The downside: it is hand-written CUDA, and it only handles dense and causal
attention. Every new attention shape needs a brand-new hand-written kernel.}

%% =============================================================
\section{The Real Problem}
%% =============================================================

This is where the project's motivation crystallises.

\tip{Researchers keep inventing new attention shapes. Mistral uses a
sliding window so each word only looks at the last 256 words. BLOOM
uses ALiBi, which adds a penalty for far-away words. LLaMA uses RoPE,
which rotates the vectors based on their position. BigBird uses block
sparse patterns. PyTorch's official fast path only knows about dense
and causal --- everything else silently falls back to the slow N-squared
code, which is sometimes hundreds of times slower. So researchers face
a brutal choice: write CUDA for a few weeks, or accept terrible
performance. There is nothing in the middle. That gap is exactly what
AttnFuse fills.}

%% =============================================================
\section{AttnFuse — The Idea}
%% =============================================================

Now the elevator pitch in concrete terms.

\tip{AttnFuse is a small language embedded inside Python. You describe
your attention shape as a short function using building blocks like
\texttt{causal}, \texttt{sliding\_window}, \texttt{alibi}, \texttt{rope},
and \texttt{softmax}. Our compiler reads that description, figures out
which steps fuse together, picks good GPU tile sizes, and generates one
fused GPU program for the whole thing. The user never writes CUDA or
Triton. The result runs at roughly the same speed as the hand-written
FlashAttention for the cases FlashAttention covers, and dozens of times
faster than PyTorch for the cases it doesn't.}

%% =============================================================
\section{What the User Writes (the Whole Thing!)}
%% =============================================================

Walk the audience through the five lines on the slide.

\tip{Here is what the user actually writes. They decorate a function with
\texttt{@af.attention}. Inside, they call \texttt{scaled\_dot\_product} on
Q and K to get the scores, then \texttt{causal} to mask out future words,
then \texttt{alibi} to add the distance penalty, then \texttt{softmax} to
turn scores into percentages, and finally they multiply by V. That is the
whole thing. The first time this function gets called, AttnFuse traces it,
compiles a fused Triton kernel, and caches it. Every later call is just a
GPU launch --- under 200 microseconds of overhead.}

%% =============================================================
\section{The Building Blocks (Combinators)}
%% =============================================================

Briefly enumerate the seven combinators. The audience should walk away
believing this is genuinely composable.

\tip{We have seven building blocks, grouped into four categories.
\textit{Scoring}: plain scaled dot product, or a fused RoPE version that
rotates the vectors as it scores them. \textit{Masking}: causal, which
hides future words, and sliding-window, which only looks back a fixed
distance. \textit{Bias}: ALiBi, the built-in distance penalty, or
additive bias, where you pass in any bias tensor you like.
\textit{Normalisation}: standard softmax, or ReLU attention for
researchers experimenting with simpler activations. You can snap these
together in any combination, and the compiler handles the rest.}

%% =============================================================
\section{How the Compiler Works}
%% =============================================================

Walk through the four-pass pipeline. Don't get academic about IR design;
just say what each pass does.

\tip{Behind the scenes there are four passes. \textit{Trace} runs the
user's Python function once with fake symbolic tensors, just to capture
the structure --- nothing actually runs on the GPU yet. \textit{Fuse}
recognises that score, mask, bias, softmax, and the final blend all
belong in one loop. \textit{Tile} picks the block sizes for the GPU,
using a small hardware-specific lookup table we tuned on the RTX 3090.
\textit{Lower} flattens the graph into a tile-loop record.
\textit{Codegen} emits a single Triton source file. Triton's just-in-time
compiler then turns that into the actual GPU machine code. The whole
thing gets cached by a hash of the graph, so reusing the same shape
costs nothing on the second call.}

%% =============================================================
\section{One Kernel, Many Variants}
%% =============================================================

This is the design trick that makes the project tractable. Emphasise
``zero runtime branching.''

\tip{Here is the design choice I am most proud of. Instead of generating
one fresh kernel per variant, we wrote one big templated kernel that
carries compile-time switches: which mask, which bias, which normaliser,
RoPE on or off. Triton's JIT is smart enough to bake those switches in
at compile time and delete the dead code paths. So at runtime, there is
not a single branch inside the hot loop --- the GPU just sees the right
specialised binary for whichever combination you asked for. Adding a new
combinator is a few lines of code-generation, not a new kernel.}

%% =============================================================
\section{The Magic Inside: Online Softmax}
%% =============================================================

Explain online softmax in words. Don't write the equations on the board
unless asked.

\tip{This is the trick that lets us reproduce FlashAttention's speed.
The naive softmax wants the whole row of scores in front of it before
it can normalise. With four thousand columns that is impossible to fit
on the GPU's tiny fast memory. The online version cheats: it walks the
row tile by tile, keeping just two running numbers per row --- the
biggest score seen so far, and a running sum. When a new tile brings a
bigger score, we rescale the old partial answer to match. At the end
of the loop you get the exact same numbers as the naive version --- not
an approximation --- but you never store the giant scoreboard. Memory
goes from N-squared to N.}

%% =============================================================
\section{Our Novel Trick — Fused RoPE}
%% =============================================================

Emphasise that this is genuinely new --- not just an implementation of
something published.

\tip{RoPE, or Rotary Position Embedding, is the way LLaMA and Mistral
encode the position of each word. Conceptually you rotate Q and K by
an angle that depends on where the word sits in the sentence. The way
everyone does this today is to rotate Q and K in two separate small
GPU programs, write the rotated versions back to slow memory, and
then run the attention kernel on top. We fold the rotation into the
attention kernel itself: rotate Q once before the inner loop, rotate
each tile of K right after we load it. Everything stays in fast
registers --- no extra trips to slow memory. The result is between
eighteen percent and two-and-a-third times faster, depending on the
sequence length. As far as I can tell, this is the first
compiler-generated fused RoPE; the few existing fused implementations
are hand-written.}

%% =============================================================
\section{Headline Results}
%% =============================================================

This is the slide that decides whether the project ``looks good.''
Read the numbers exactly and call out the punchline.

\tip{Now the numbers, all measured on an RTX 3090 with sixteen-bit floats
at a four thousand ninety six token sequence length. For dense attention
we get fifty-nine teraflops, PyTorch's SDPA gets seventy --- about
eighty-five percent of the hand-tuned CUDA. For causal it is one hundred
and five versus one hundred and twenty-two --- eighty-six percent. So
for the shapes PyTorch already accelerates, we are competitive. The
interesting numbers are the next two. For sliding-window we get two
hundred and ninety-eight teraflops versus PyTorch's seven point eight
--- a thirty-eight times speedup. For causal-plus-ALiBi we get ninety-eight
versus five point six --- seventeen times speedup. AttnFuse is the
only sub-quadratic option for these shapes without writing custom CUDA.}

%% =============================================================
\section{Memory + Multi-Dtype}
%% =============================================================

Hit the memory savings and the dtype coverage together.

\tip{The other big win is memory. At four thousand tokens, the naive
PyTorch attention allocates three point two gigabytes of GPU memory
just for the scoreboard; AttnFuse uses two hundred megabytes --- about
thirty-two times less. This is the same memory savings FlashAttention
gives you, but you get it automatically for any user-written variant.
We also run on three precisions: sixteen-bit float, sixteen-bit brain
float, and full thirty-two-bit float, each with its own tile config.
In thirty-two bit, PyTorch loses its flash path and falls back to
slow code; AttnFuse stays fast and ends up three-point-two times
faster than SDPA in that regime. All of this is verified by twenty-five
GPU unit tests across every variant.}

%% =============================================================
\section{Costs, Limitations, and What We Built}
%% =============================================================

Honest closing slide. Be upfront about what's missing; then the summary.

\tip{Nothing is free. The first time you call a new attention shape,
the compile takes about a second and a half. After that the compiled
kernel is cached forever --- subsequent calls are under two hundred
microseconds. We do only the forward pass right now; the backward pass
for training is future work. We are single-GPU; we don't support block
sparse patterns yet, though the IR has a slot for them. To summarise:
seven combinators, a two-level IR, a four-pass compiler, the fused RoPE
kernel, twenty-five GPU tests, three precisions, four sequence lengths,
three baselines, and full evaluation. The bottom line is a declarative
DSL that matches FlashAttention 2 within fifteen percent on the cases
PyTorch already accelerates, and beats it seventeen to thirty-eight
times on the cases it doesn't. Thank you --- happy to take questions.}

%% =============================================================
\section*{Appendix: Likely Questions \& Short Answers}
%% =============================================================

\paragraph{Q. Why Triton, not CUDA?}
Triton is a higher-level language than CUDA. We can generate it from a
template without dealing with PTX assembly or shared-memory bank
conflicts. The cost is roughly fifteen percent versus hand-tuned CUDA,
which is acceptable for kernel generality.

\paragraph{Q. What is the memory cost of online softmax?}
Per row, only the running max and the running sum --- two numbers in
registers. Nothing extra in HBM. That is what reduces $O(N^2)$ memory
to $O(N)$.

\paragraph{Q. How is the cache key computed?}
A SHA-1 hash of the high-level graph's structural signature, including
the rope flag, mask kind, bias kind, head dim, and dtype. Two graphs
that hash the same share a compiled kernel.

\paragraph{Q. What about training (backward)?}
Not implemented. The backward pass is a separate kernel with different
tile scheduling, as in FlashAttention-2. Adding it would be the
natural next milestone.

\paragraph{Q. Did you profile with Nsight Compute?}
\texttt{ncu} was not available on the evaluation host. The TFLOPS and
HBM bandwidth in the roofline are computed analytically from measured
latency with the correct FLOP and traffic formulas per variant. We note
this in the report.

\paragraph{Q. Is this faster than xFormers or Liger Kernel?}
For dense and causal, those libraries match SDPA. AttnFuse's edge is
covering combinations they don't ship --- sliding-window, ALiBi,
fused RoPE, custom additive bias --- without writing a new kernel per
combination.

\end{document}
